\documentclass[10pt,a4paper]{amsart}
\usepackage{amsmath,amssymb,amsthm,graphicx}
% Retained preamble for the Stacks Project chapter snapshot in this directory.
% The chapter's own preamble.tex is a copy of the shared Stacks Project preamble; the block
% below is sliced MECHANICALLY, line for line, from it. Nothing was added, removed or
% redefined; the wrapper class, the xr-hyper/hyperref wiring and the externaldocument list
% are omitted because the assembler supplies the document class.
%
% lines taken from preamble.tex: 3, 5, 9, 11, 12, 156, 157, 158, ... (51 lines total)
%
% package coverage restored from the chapter's own preamble (2026-10-13): xy, xy-2cell

\usepackage[all]{xy}
\xyoption{2cell}
\UseAllTwocells

\newenvironment{reference}{\comment}{\endcomment}
\newenvironment{slogan}{\comment}{\endcomment}
\newenvironment{history}{\comment}{\endcomment}
\theoremstyle{plain}
\newtheorem{theorem}[subsection]{Theorem}
\newtheorem{proposition}[subsection]{Proposition}
\newtheorem{lemma}[subsection]{Lemma}
\theoremstyle{definition}
\newtheorem{definition}[subsection]{Definition}
\newtheorem{example}[subsection]{Example}
\newtheorem{exercise}[subsection]{Exercise}
\newtheorem{situation}[subsection]{Situation}
\theoremstyle{remark}
\newtheorem{remark}[subsection]{Remark}
\newtheorem{remarks}[subsection]{Remarks}
\numberwithin{equation}{subsection}
\def\lim{\mathop{\mathrm{lim}}\nolimits}
\def\colim{\mathop{\mathrm{colim}}\nolimits}
\def\Spec{\mathop{\mathrm{Spec}}}
\def\Hom{\mathop{\mathrm{Hom}}\nolimits}
\def\Ext{\mathop{\mathrm{Ext}}\nolimits}
\def\SheafHom{\mathop{\mathcal{H}\!\mathit{om}}\nolimits}
\def\SheafExt{\mathop{\mathcal{E}\!\mathit{xt}}\nolimits}
\def\Sch{\mathit{Sch}}
\def\Mor{\mathop{\mathrm{Mor}}\nolimits}
\def\Ob{\mathop{\mathrm{Ob}}\nolimits}
\def\Sh{\mathop{\mathit{Sh}}\nolimits}
\def\NL{\mathop{N\!L}\nolimits}
\def\CH{\mathop{\mathrm{CH}}\nolimits}
\def\proetale{{pro\text{-}\acute{e}tale}}
\def\etale{{\acute{e}tale}}
\def\QCoh{\mathit{QCoh}}
\def\Ker{\mathop{\mathrm{Ker}}}
\def\Im{\mathop{\mathrm{Im}}}
\def\Coker{\mathop{\mathrm{Coker}}}
\def\Coim{\mathop{\mathrm{Coim}}}
\DeclareMathSymbol{\boxtimes}{\mathbin}{AMSa}{"02}
\def\QCohstack{\mathcal{QC}\!\mathit{oh}}
\def\Cohstack{\mathcal{C}\!\mathit{oh}}
\def\Spacesstack{\mathcal{S}\!\mathit{paces}}
\def\Quotfunctor{\mathrm{Quot}}
\def\Hilbfunctor{\mathrm{Hilb}}
\def\Curvesstack{\mathcal{C}\!\mathit{urves}}
\def\Polarizedstack{\mathcal{P}\!\mathit{olarized}}
\def\Complexesstack{\mathcal{C}\!\mathit{omplexes}}
\def\Pic{\mathop{\mathrm{Pic}}\nolimits}
\def\Picardstack{\mathcal{P}\!\mathit{ic}}
\def\Picardfunctor{\mathrm{Pic}}
\def\Deformationcategory{\mathcal{D}\!\mathit{ef}}
\title{Three equivalent forms of regularity for a sequence that does not generate the unit ideal}
\author{The Stacks Project Authors}
\date{stacks-project:algebra.tex:local-snapshot-20261012; GNU FDL 1.2 or later}
\begin{document}
\maketitle
\noindent\textit{Editorial note.} All delivered mathematical text is the author's own, copied from the pinned Stacks Project chapter file algebra.tex; no statement or proof step was rewritten, reordered or removed.\par\medskip
\noindent\textit{Editorial note.} Preamble provenance, as required for this source: the frozen snapshot stages the chapter's own preamble file beside it (retained-preamble.tex, 63 lines, 8116 bytes, the same shared Stacks-wide preamble the published chow and etale-cohomology cards use); the first line of the pinned chapter file is \textbackslash{}input\{preamble\}, so no preamble text lives inside the chapter itself. The card takes that file's lines 8-18 (the verbatim package and the three \textbackslash{}comment-based environments reference/slogan/history at lines 9-12, section 13.2, together with the Xy-pic setup) and its lines 154 to the runtime end of the file (the native theorem declarations and mathematical macros). Not one character was changed, no package was added, no macro was defined, and the excerpt serves only as preamble, never as a span coordinate.\par\medskip
\noindent\textit{Editorial note.} The same-chapter pointers are located at their \textbackslash{}label lines inside the pinned chapter file; the cross-chapter pointers are located in the project's own tag list staged in the snapshot (sources/stacks-chow/tags.txt), at the verbatim lines carrying their canonical tags.\par\medskip
\noindent\textit{Editorial note.} The counted spans contain no \textbackslash{}ref and no \textbackslash{}cite, so the card carries no navigation stub and no bibliography.\par\medskip
\noindent\textit{Allowed standard prerequisites (editorial).}\begin{itemize}
\item Kernels of multiplication on quotients. Used at the start of the counted proof (lines 17028-17031). Turns the two-element case into a comparison of kernels.
\item Regularity inherited by subsequences. Used in the induction step. Lets the induction on the length be applied to shorter sequences.
\end{itemize}\medskip
\begin{lemma}
\label{lemma-regular-sequence-in-polynomial-ring}
Let $R$ be a ring. Let $f_1, \ldots, f_r \in R$ which do not generate
the unit ideal. The following are equivalent:
\begin{enumerate}
\item any permutation of $f_1, \ldots, f_r$ is a regular sequence,
\item any subsequence of $f_1, \ldots, f_r$ (in the given order) is
a regular sequence, and
\item $f_1x_1, \ldots, f_rx_r$ is a regular sequence in the polynomial
ring $R[x_1, \ldots, x_r]$.
\end{enumerate}
\end{lemma}

\begin{proof}
It is clear that (1) implies (2). We prove (2) implies (1) by induction
on $r$. The case $r = 1$ is trivial. The case $r = 2$ says that if
$a, b \in R$ are a regular sequence and $b$ is a nonzerodivisor, then
$b, a$ is a regular sequence. This is clear because the kernel of
$a : R/(b) \to R/(b)$ is isomorphic to the kernel of $b : R/(a) \to R/(a)$
if both $a$ and $b$ are nonzerodivisors. The case $r > 2$. Assume
(2) holds and say we want to prove $f_{\sigma(1)}, \ldots, f_{\sigma(r)}$
is a regular sequence for some permutation $\sigma$. We already know
that $f_{\sigma(1)}, \ldots, f_{\sigma(r - 1)}$ is a regular sequence
by induction. Hence it suffices to show that $f_s$ where $s = \sigma(r)$
is a nonzerodivisor modulo $f_1, \ldots, \hat f_s, \ldots, f_r$.
If $s = r$ we are done. If $s < r$, then note that $f_s$ and $f_r$
are both nonzerodivisors in the ring
$R/(f_1, \ldots, \hat f_s, \ldots, f_{r - 1})$
(by induction hypothesis again). Since we know $f_s, f_r$ is a
regular sequence in that ring we conclude by the case of sequence of length
$2$ that $f_r, f_s$ is too.

\medskip\noindent
Note that $R[x_1, \ldots, x_r]/(f_1x_1, \ldots, f_ix_i)$ as an $R$-module
is a direct sum of the modules
$$
R/I_E \cdot x_1^{e_1} \ldots x_r^{e_r}
$$
indexed by multi-indices $E = (e_1, \ldots, e_r)$ where
$I_E$ is the ideal generated by $f_j$ for $1 \leq j \leq i$
with $e_j > 0$. Hence $f_{i + 1}x_i$ is a nonzerodivisor on this if
and only if $f_{i + 1}$ is a nonzerodivisor on $R/I_E$ for all $E$.
Taking $E$ with all positive entries, we see that $f_{i + 1}$
is a nonzerodivisor on $R/(f_1, \ldots, f_i)$. Thus (3) implies (2).
Conversely, if (2) holds, then any subsequence of
$f_1, \ldots, f_i, f_{i + 1}$ is a regular sequence
in particular $f_{i + 1}$ is a nonzerodivisor on all $R/I_E$.
In this way we see that (2) implies (3).
\end{proof}

\end{document}
